\documentclass[11pt,a4paper]{article}
\usepackage{curso}
\begin{document}
\portada{Ejercicios --- Lección}{1.5}{El principio de Landauer}
Constantes: $\kB=1{,}380649\times10^{-23}$ J/K, $\ln2\approx0{,}6931$, $1\,\text{eV}=1{,}602\times10^{-19}$ J.

\section{Ejercicios básicos}
\begin{enumerate}[label=\textbf{B\arabic*.}]
\item Calcule $\kB T\ln2$ en julios y en eV a $T=300$ K, $T=77$ K y $T=4$ K.
\item ¿Cuál es el calor mínimo para borrar 1 GB ($8\times10^9$ bits) a 300 K?
\item Un bit tiene $p(0)=0{,}9$. ¿Cuál es el costo mínimo para reiniciarlo a 0 a 300 K?
\item Indique si cada función es lógicamente reversible: NOT, AND, XOR de dos bits, $(a,b)\mapsto(a,a\oplus b)$ (CNOT).
\end{enumerate}

\section{Ejercicios intermedios}
\begin{enumerate}[label=\textbf{I\arabic*.}]
\item Derive el trabajo de compresión isotérmica de un gas de una partícula de $V$ a $V/2$ y explique por qué $Q=W$.
\item Con entradas uniformes e independientes, calcule la información perdida por AND, OR y XOR; ¿cuál es su costo mínimo de Landauer?
\item Un chip realiza $10^{18}$ borrados de bit por segundo a 300 K. ¿Qué potencia mínima impone Landauer? Compárela con 1 W.
\item Usando la tabla de simulación de las notas, ajuste $\langle W\rangle=\kB T\ln2+a/\tau$ con los puntos $\tau\ge10$ y estime $a$.
\end{enumerate}

\section{Ejercicios avanzados}
\begin{enumerate}[label=\textbf{A\arabic*.}]
\item Para el potencial $V=E_b(x^4-2x^2)$, use la fórmula de Kramers $k\approx\frac{\mu\sqrt{V''(x_{\min})|V''(x_{\max})|}}{2\pi}e^{-E_b/\kB T}$ para estimar la vida media de un bit con $E_b=8\kB T$ y $E_b=40\kB T$ (unidades $\mu=\kB T=1$).
\item Si el borrado falla con probabilidad $\epsilon$ (el bit queda en 1), muestre que la cota se convierte en $Q\ge\kB T\ln2\,[1-h(\epsilon)]$. Evalúe para $\epsilon=0{,}01$.
\item Explique, con la descomposición calcular--copiar--descalcular de Bennett, cuántos bits deben borrarse al final de un sumador completo reversible que conserva solo la suma y el acarreo.
\end{enumerate}

\section{Soluciones}
\begin{solucion}
\textbf{B1.} $300$ K: $2{,}871\times10^{-21}$ J $=0{,}0179$ eV. $77$ K: $7{,}37\times10^{-22}$ J $=0{,}0046$ eV. $4$ K: $3{,}83\times10^{-23}$ J $=2{,}39\times10^{-4}$ eV.

\textbf{B2.} $8\times10^9\times2{,}871\times10^{-21}=2{,}30\times10^{-11}$ J.

\textbf{B3.} $h(0{,}9)=-0{,}9\log_20{,}9-0{,}1\log_20{,}1=0{,}469$ bits. $Q_{\min}=0{,}469\cdot2{,}871\times10^{-21}=1{,}35\times10^{-21}$ J.

\textbf{B4.} NOT: sí (biyección). AND: no. XOR de dos entradas a una salida: no (2 bits $\to$1 bit). CNOT: sí (es su propia inversa).

\textbf{I1.} $W=-\int_V^{V/2}\frac{\kB T}{V'}\dd V'=\kB T\ln2$. Para un gas ideal $U$ depende solo de $T$; a $T$ constante $\Delta U=0$ y la primera ley da $Q=W$.

\textbf{I2.} AND: salida con $P(1)=1/4$, $H=0{,}811$; pérdida $1{,}189$ bits. OR: $P(1)=3/4$, misma pérdida $1{,}189$ bits. XOR: $P(1)=1/2$, $H=1$; pérdida 1 bit. Costos: $1{,}189\,\kB T\ln2$, $1{,}189\,\kB T\ln2$ y $\kB T\ln2$.

\textbf{I3.} $P=10^{18}\times2{,}871\times10^{-21}=2{,}87\times10^{-3}$ W, unas 350 veces menos que 1 W.

\textbf{I4.} Con $x=1/\tau$ y $y=\langle W\rangle-\ln2$: puntos $(0{,}1;0{,}42)$, $(0{,}05;0{,}21)$, $(0{,}02;0{,}09)$, $(0{,}01;0{,}04)$. Mínimos cuadrados por el origen: $a=\sum xy/\sum x^2\approx(0{,}042+0{,}0105+0{,}0018+0{,}0004)/(0{,}01+0{,}0025+0{,}0004+0{,}0001)\approx4{,}2$.

\textbf{A1.} $V''(\pm1)=8E_b$, $|V''(0)|=4E_b$; prefactor $\sqrt{32}E_b/2\pi\approx0{,}90E_b$. Con $E_b=8$: $k\approx7{,}2\,e^{-8}\approx2{,}4\times10^{-3}$, vida media $\sim400$ unidades de tiempo. Con $E_b=40$: $k\approx36\,e^{-40}\approx1{,}5\times10^{-16}$, vida media $\sim7\times10^{15}$. La memoria pasa de volátil a prácticamente permanente: la confiabilidad exige barreras de decenas de $\kB T$.

\textbf{A2.} La entropía final de la memoria es $\kB\ln2\,h(\epsilon)$, así que $\Delta S_{\text{mem}}=-\kB\ln2\,[1-h(\epsilon)]$ y la segunda ley da la cota. Para $\epsilon=0{,}01$: $h=0{,}0808$, $Q\ge0{,}919\,\kB T\ln2$.

\textbf{A3.} El sumador completo mapea 3 bits $(a,b,c_{\text{in}})$ a 2 bits (suma, acarreo): no es inyectivo. En la versión de Bennett se calcula de forma reversible guardando una entrada (p. ej. $a$) como basura, se copian suma y acarreo, y se descalcula. Si al final se quiere conservar solo la salida, debe borrarse la información que falta para reconstruir la entrada: en promedio $3-H(\text{suma},\text{acarreo})=3-1{,}811=1{,}19$ bits con entradas uniformes (las salidas $(s,c)$ tienen probabilidades $1/8,3/8,3/8,1/8$).
\end{solucion}

\section{Ejercicios adicionales (sin solución)}
\begin{enumerate}[label=\textbf{P\arabic*.}]
\item Una memoria de $k$ niveles (pdit) uniformes se reinicia. ¿Cuál es su cota de Landauer?
\item Diseñe un protocolo de borrado en el pozo doble que no use inclinación sino un desplazamiento de un pozo. ¿Cambia la cota?
\item Implemente en Python la simulación del pozo doble y reproduzca la tabla $\langle W\rangle(\tau)$ (laboratorio L2).
\item ¿Por qué la máquina de Szilard no viola la segunda ley? Escriba el balance completo del ciclo.
\end{enumerate}
\end{document}
